% Options for packages loaded elsewhere
\PassOptionsToPackage{unicode}{hyperref}
\PassOptionsToPackage{hyphens}{url}
\documentclass[
  11pt,
]{article}
\usepackage{xcolor}
\usepackage[left=3cm,right=3cm,top=2.5cm,bottom=2.5cm]{geometry}
\usepackage{amsmath,amssymb}
\setcounter{secnumdepth}{5}
\usepackage{iftex}
\ifPDFTeX
  \usepackage[T1]{fontenc}
  \usepackage[utf8]{inputenc}
  \usepackage{textcomp} % provide euro and other symbols
\else % if luatex or xetex
  \usepackage{unicode-math} % this also loads fontspec
  \defaultfontfeatures{Scale=MatchLowercase}
  \defaultfontfeatures[\rmfamily]{Ligatures=TeX,Scale=1}
\fi
\usepackage{lmodern}
\ifPDFTeX\else
  % xetex/luatex font selection
\fi
% Use upquote if available, for straight quotes in verbatim environments
\IfFileExists{upquote.sty}{\usepackage{upquote}}{}
\IfFileExists{microtype.sty}{% use microtype if available
  \usepackage[]{microtype}
  \UseMicrotypeSet[protrusion]{basicmath} % disable protrusion for tt fonts
}{}
\usepackage{setspace}
\makeatletter
\@ifundefined{KOMAClassName}{% if non-KOMA class
  \IfFileExists{parskip.sty}{%
    \usepackage{parskip}
  }{% else
    \setlength{\parindent}{0pt}
    \setlength{\parskip}{6pt plus 2pt minus 1pt}}
}{% if KOMA class
  \KOMAoptions{parskip=half}}
\makeatother
\usepackage{longtable,booktabs,array}
\usepackage{caption}
\captionsetup[table]{skip=6pt}
\newcounter{none} % for unnumbered tables
\usepackage{calc} % for calculating minipage widths
% Correct order of tables after \paragraph or \subparagraph
\usepackage{etoolbox}
\makeatletter
\patchcmd\longtable{\par}{\if@noskipsec\mbox{}\fi\par}{}{}
\makeatother
% Allow footnotes in longtable head/foot
\IfFileExists{footnotehyper.sty}{\usepackage{footnotehyper}}{\usepackage{footnote}}
\makesavenoteenv{longtable}
\usepackage{graphicx}
\makeatletter
\newsavebox\pandoc@box
\newcommand*\pandocbounded[1]{% scales image to fit in text height/width
  \sbox\pandoc@box{#1}%
  \Gscale@div\@tempa{\textheight}{\dimexpr\ht\pandoc@box+\dp\pandoc@box\relax}%
  \Gscale@div\@tempb{\linewidth}{\wd\pandoc@box}%
  \ifdim\@tempb\p@<\@tempa\p@\let\@tempa\@tempb\fi% select the smaller of both
  \ifdim\@tempa\p@<\p@\scalebox{\@tempa}{\usebox\pandoc@box}%
  \else\usebox{\pandoc@box}%
  \fi%
}
% Set default figure placement to htbp
\def\fps@figure{htbp}
\makeatother
\setlength{\emergencystretch}{3em} % prevent overfull lines
\providecommand{\tightlist}{%
  \setlength{\itemsep}{0pt}\setlength{\parskip}{0pt}}
%% tools/stamp.tex
%% Provenance footer for PDFs rendered from this project.
%% Vendored by zzcollab; this file is not project-specific. The
%% footer prints the source document and the time of rendering.
%% The git version is defined here too, and so is available to a
%% project that wants it on the page, but it is left off the
%% default footer: it already names the staged copy in share/ and
%% appears in its manifest row. All values are supplied at render
%% time by tools/stamp-render.R, which writes a small generated
%% file that redefines the macros below. The \providecommand
%% defaults keep a bare LaTeX compile working when the document is
%% built without the wrapper.
\usepackage{fancyhdr}
\usepackage{xcolor}
\providecommand{\stampsource}{\detokenize{(source not recorded)}}
\providecommand{\stamptime}{(time not recorded)}
\providecommand{\stampversion}{\detokenize{(version not recorded)}}
\setlength{\footskip}{40pt}
\newcommand{\stampfootcontent}{%
  \footnotesize\color{gray}%
  \parbox[b]{\textwidth}{%
    \stamptime\hfill\thepage\\[2pt]%
    \ttfamily\raggedright\stampsource}}
\pagestyle{fancy}
\fancyhf{}
\fancyfoot[C]{\stampfootcontent}
\renewcommand{\headrulewidth}{0pt}
\renewcommand{\footrulewidth}{0.4pt}
%% The title page uses the `plain` page style; redefine it so the
%% stamp appears there too.
\fancypagestyle{plain}{%
  \fancyhf{}%
  \fancyfoot[C]{\stampfootcontent}%
  \renewcommand{\headrulewidth}{0pt}%
  \renewcommand{\footrulewidth}{0.4pt}}
\renewcommand{\stampsource}{\textasciitilde{}/\allowbreak{}prj/\allowbreak{}res/\allowbreak{}04-runin-power-analysis/\allowbreak{}runinpower/\allowbreak{}analysis/\allowbreak{}report/\allowbreak{}02-correlation/\allowbreak{}bullets.Rmd}
\renewcommand{\stamptime}{2026-08-16 16:15 PDT}
\renewcommand{\stampversion}{\detokenize{276ac7f-wip}}
\usepackage{amsmath}
\usepackage{amssymb}
\usepackage{booktabs}
\usepackage{longtable}
\usepackage{float}
\usepackage[mathlines]{lineno}
\linenumbers
\DeclareMathOperator{\Var}{Var}
\DeclareMathOperator{\Cov}{Cov}
\DeclareMathOperator{\Corr}{Corr}
\DeclareMathOperator{\tr}{tr}
\usepackage{bookmark}
\IfFileExists{xurl.sty}{\usepackage{xurl}}{} % add URL line breaks if available
\urlstyle{same}
% fallback for those not using the hyperref driver hyperxmp:
\makeatletter
\@ifundefined{xmpquote}{\newcommand{\xmpquote}[1]{#1}}{}
\makeatother
\hypersetup{
  pdftitle={Structured summary: \textbar{} Power Analysis for Run-In Clinical Trials under AR(1) and General Correlation Structures},
  pdfauthor={Ronald G. Thomas, Ph.D.},
  hidelinks,
  pdfcreator={LaTeX via pandoc}}

\title{Structured summary: \textbar{} Power Analysis for Run-In Clinical Trials under AR(1) and General Correlation Structures}
\author{Ronald G. Thomas, Ph.D.\thanks{Wertheim School of Public Health, UCSD; ORCID: 0000-0003-1686-4965}}
\date{August 16, 2026}

\begin{document}
\maketitle

{
\setcounter{tocdepth}{2}
\tableofcontents
}
\setstretch{1.1}
\emph{This document is a structured, section-by-section summary of the key points argued in the companion manuscript, ``| Power Analysis for Run-In Clinical Trials under AR(1) and General Correlation Structures.'' It is provided as a supplementary reading aid and does not stand on its own; all notation, derivations, simulation detail, and citations are in the main manuscript.}

\bigskip\noindent

\rule{\textwidth}{0.4pt}\bigskip

\section{Introduction}\label{introduction}

\textbf{The efficiency of run-in designs is governed by the correlation structure of the repeated measurements.}

\begin{itemize}\setlength{\itemsep}{2pt}\setlength{\parskip}{0pt}
\item Under compound symmetry the correlation between change scores is constant.
\item Run-in observations then inform the random slope through simple averaging.
\item Under AR(1) the correlation decays with separation, so run-in information depends on distance from the treatment period.
\end{itemize}

\textbf{This paper extends prior compound-symmetric run-in frameworks to AR(1) residual correlation.}

\begin{itemize}\setlength{\itemsep}{2pt}\setlength{\parskip}{0pt}
\item Paper 1 assumed the rank-one-plus-identity change-score covariance from independent errors.
\item Wang (2019) extended run-in to two-period mixed models but kept compound symmetry.
\item We instead adopt AR(1) residual correlation with geometric decay in lag.
\end{itemize}

\textbf{Existing literature treats autocorrelation and change-score variation but not the run-in setting.}

\begin{itemize}\setlength{\itemsep}{2pt}\setlength{\parskip}{0pt}
\item Winkens (2007) studied optimal designs for autocorrelated data outside the run-in context.
\item Diggle and Verbeke give general treatments of longitudinal correlation structures.
\item Munoz et al. (1992) embedded compound symmetry and AR(1) in a single damped-exponential family, the natural generalization of the two structures contrasted here.
\item Chambless (1993) addressed intra-individual variation in change scores, a concern amplified under AR(1).
\end{itemize}

\textbf{The closed-form AR(1) derivation covers a design case the reference longpower toolkit does not.}

\begin{itemize}\setlength{\itemsep}{2pt}\setlength{\parskip}{0pt}
\item Its slope-based closed-form routines assume independent or compound-symmetric residuals.
\item Its \texttt{power.mmrm.ar1()} routine implements the analytic AR(1) formula of Lu, Luo, and Chen (2008), but without run-in or common close observations under a random slope.
\item Our formula supplies that case and collapses to the compound-symmetric structure of Paper 1 as $\rho \to 0$.
\end{itemize}

\subsection{Marginal covariance}\label{marginal-covariance}

\textbf{The AR(1) change-score covariance loses rank-one structure but the Woodbury reduction survives.}

\begin{itemize}\setlength{\itemsep}{2pt}\setlength{\parskip}{0pt}
\item The matrix $R$ is no longer rank-one-plus-identity, so Sherman-Morrison does not apply.
\item The Woodbury identity still reduces the inverse covariance to a scalar inversion.
\item This holds because the random-effects covariance $G$ remains scalar.
\end{itemize}

\subsection{Woodbury computation}\label{woodbury-computation}

\textbf{The Woodbury scalar term persists, but the residual inverse must now be computed differently.}

\begin{itemize}\setlength{\itemsep}{2pt}\setlength{\parskip}{0pt}
\item The correction term $H$ remains a scalar as in Paper 1.
\item The residual inverse $R^{-1}$ must be obtained numerically.
\item Equivalently it follows from the AR(1) tridiagonal inverse rather than Sherman-Morrison.
\end{itemize}

\subsection{Cross-covariance between phase means}\label{cross-covariance-between-phase-means}

\textbf{A moment-based alternative derives the pre-post variance directly from the phase means.}

\begin{itemize}\setlength{\itemsep}{2pt}\setlength{\parskip}{0pt}
\item It computes the variance of the test statistic from the moments of the phase means.
\item The run-in mean averages the run-in observations.
\item The treatment-period mean averages the treatment-period observations.
\item The construction follows the summary-measures approach of Frison and Pocock (1992).
\end{itemize}

\section{Relationship between the two approaches}\label{relationship-between-the-two-approaches}

\textbf{The GLS and moment-based approaches target distinct variance quantities.}

\begin{itemize}\setlength{\itemsep}{2pt}\setlength{\parskip}{0pt}
\item The GLS variance is the minimum-variance bound under the assumed model.
\item The moment-based variance is that of the simple averaged change-score estimator.
\item The two therefore answer different efficiency questions.
\end{itemize}

\textbf{The two approaches do not coincide even without a random slope.}

\begin{itemize}\setlength{\itemsep}{2pt}\setlength{\parskip}{0pt}
\item At $\sigma_b^2 = 0$ the GLS estimator still weights observations through the AR(1) residual covariance and targets the slope difference.
\item The moment-based quantity is the variance of the unweighted phase-mean difference, a different estimand.
\item Numerically, at $J_0 = 2$, $J_1 = 2$, $\rho = 0.5$, $\sigma^2 = 1$, $\sigma_b^2 = 0$: GLS 0.636 versus moment-based 1.219.
\end{itemize}

\textbf{The moment-based computation serves the simple pre-post estimator, not as a bound on the GLS variance.}

\begin{itemize}\setlength{\itemsep}{2pt}\setlength{\parskip}{0pt}
\item It requires no specification of the random-slope variance.
\item It is the exact variance of the unweighted pre-post difference an analyst might report without fitting a mixed model.
\item Because the estimands differ, its ratio to the GLS variance is descriptive, not an efficiency bound.
\end{itemize}

\section{Discussion}\label{discussion}

\textbf{The autocorrelation structure matters for run-in designs only when the autocorrelation is weak relative to the random-slope signal.}

\begin{itemize}\setlength{\itemsep}{2pt}\setlength{\parskip}{0pt}
\item When the slope-to-noise ratio is large, the random-slope term dominates the change-score covariance and the residual correlation structure is second order.
\item The AR(1) and compound-symmetry variances then differ negligibly, and the simpler Paper 1 formula suffices.
\item The two structures diverge most when the slope-to-noise ratio is small and run-in observations are distant from randomization.
\end{itemize}

\textbf{At a matched nominal $\rho$, compound symmetry is generally conservative, with a narrow anticonservative corner.}

\begin{itemize}\setlength{\itemsep}{2pt}\setlength{\parskip}{0pt}
\item Positive serial correlation reduces the residual variance of every change score, so the AR(1) variance falls below the compound-symmetric value for most designs.
\item The discrepancy grows with $\rho$; near $\rho = 0.95$ the AR(1) variance can fall below one tenth of the compound-symmetric value.
\item The exception is large $J_0$ at small $\rho$, where distant run-in observations lose information and compound symmetry understates the variance by at most about 6\%.
\end{itemize}

\textbf{The numerical comparisons give a practical decision rule for when the simpler formula is safe.}

\begin{itemize}\setlength{\itemsep}{2pt}\setlength{\parskip}{0pt}
\item At a matched nominal $\rho$, the compound-symmetry formula of Paper 1 is a conservative and convenient default for most designs.
\item With many run-in observations and low expected autocorrelation, the AR(1) formula should be used directly; the compound-symmetric formula can understate the variance by up to about 6\% there.
\item When $\rho$ is uncertain, evaluating the AR(1) formula over a plausible range of $\rho$ and taking the largest variance is a safe default.
\end{itemize}

\section{Data availability statement}\label{data-availability-statement}

\textbf{The implementing package, tests, and bibliography are openly available in the project repository.}

\begin{itemize}\setlength{\itemsep}{2pt}\setlength{\parskip}{0pt}
\item The package implements the AR(1) variance formulas.
\item Validation tests and the shared four-paper bibliography are included.
\item Specific paths appear in the Reproducibility section below.
\end{itemize}

\section{Reproducibility}\label{reproducibility}

\textbf{The manuscript was rendered with R 4.4.x using a small set of principal packages.}

\begin{itemize}\setlength{\itemsep}{2pt}\setlength{\parskip}{0pt}
\item The in-package source supplies the AR(1) machinery.
\item Testing uses the tinytest harness.
\item The manuscript build uses rmarkdown and bookdown.
\end{itemize}

\textbf{All reported quantities are deterministic.}

\begin{itemize}\setlength{\itemsep}{2pt}\setlength{\parskip}{0pt}
\item Given fixed inputs the variance formulas exercise no random-number generator.
\item No Monte Carlo computation appears in this manuscript.
\item Re-running the chunks therefore reproduces every table and figure exactly.
\end{itemize}

\textbf{All source, companion manuscripts, and bibliography reside at documented repository paths.}

\begin{itemize}\setlength{\itemsep}{2pt}\setlength{\parskip}{0pt}
\item The R package source is in the package source directory.
\item The companion manuscripts for Papers 1, 3, and 4 have their own analysis subdirectories.
\item This manuscript source and the shared bibliography reside in the correlation-paper directory.
\end{itemize}

\end{document}
